\section{协方差、相关性与条件波动}
\label{sec:05-Probability/04-Expectation-and-Moments/03}

你持有两只股票，过去一年每日收益率的统计如下：

\begin{itemize}
\item 股票 A：日均收益 0.05\%，日收益方差 4（即标准差 2\%）。
\item 股票 B：日均收益 0.08\%，日收益方差 9（即标准差 3\%）。
\end{itemize}

你把资金对半分开，各投 50\%。组合的日收益率是

\[
R=0.5R_A+0.5R_B.
\]

你想知道组合收益率的方差——它衡量这个投资的日常波动风险。

如果两只股票独立波动，方差直接按权重平方相加：

\[
\Var(R)=0.5^2\cdot4+0.5^2\cdot9=1+2.25=3.25.
\]

标准差约 1.80\%。这个数字假设了两只股票各涨各的，互不相干。但现实中当市场整体下行，它们大概率一起跌。你的直觉也对：两个资产如果总是同方向摆动，组合风险应该比独立情况更大才对。

怎么定量刻画 ``同方向摆动'' 的程度？

\subsection{协方差：共同偏离的信号}

对于两个随机变量 \(X,Y\)，各自有自己的均值 \(\mu_X,\mu_Y\)。当 \(X\) 高于均值时，\(Y\) 也倾向于高于均值吗？协方差同时检查二者的偏离方向：

\[
\Cov(X,Y)=\E[(X-\mu_X)(Y-\mu_Y)].
\]

如果 \(X\) 和 \(Y\) 倾向于同方向偏离各自的均值，乘积 \((X-\mu_X)(Y-\mu_Y)\) 平均为正；反方向偏离则平均为负；无系统关联则积的正负互相抵消，接近零。分开来看：

\begin{itemize}
\item \(\Cov>0\)：往往同涨同跌。
\item \(\Cov<0\)：往往此涨彼跌——对冲关系。
\item \(\Cov=0\)：没有线性共同变化（但完全可能有非线性依赖——后面会举例）。
\end{itemize}

展开定义式可以拿到一个更常用的计算公式：

\[
\Cov(X,Y)=\E[XY]-\E[X]\E[Y].
\]

验证：把 \((X-\mu_X)(Y-\mu_Y)=XY-\mu_X Y-\mu_Y X+\mu_X\mu_Y\) 取期望，中间两项分别给出 \(-\mu_X\mu_Y\) 和 \(-\mu_Y\mu_X\)，加上最后的 \(+\mu_X\mu_Y\)，凑成这个减法形式。

方差是协方差的特例：

\[
\Var(X)=\Cov(X,X).
\]

\subsection{和的方差：为什么多出交叉项}

现在回到两只股票的组合。你不应该盲目地把方差相加——展开和的方差：

\[
\begin{aligned}
\Var(X+Y)&=\E[(X+Y-(\mu_X+\mu_Y))^2]\\
&=\E[((X-\mu_X)+(Y-\mu_Y))^2]\\
&=\E[(X-\mu_X)^2]+\E[(Y-\mu_Y)^2]+2\E[(X-\mu_X)(Y-\mu_Y)]\\
&=\Var(X)+\Var(Y)+2\Cov(X,Y).
\end{aligned}
\]

交叉项 \(2\Cov(X,Y)\) 就是依赖性带来的额外波动贡献。如果两只股票正相关，交叉项为正，组合方差大于各自方差之和；如果它们负相关，交叉项为负，组合方差反而小于各自方差之和——这就是对冲降低风险的数学根基。

对于权重情况：

\[
\Var(0.5X+0.5Y)=0.25\Var(X)+0.25\Var(Y)+0.5\Cov(X,Y).
\]

用一组具体数值感受幅度。假设 \(\Cov(A,B)=4.5\)（意味着相关性不弱），那么

\[
\Var(R)=0.25\times4+0.25\times9+0.5\times4.5=1+2.25+2.25=5.5.
\]

标准差约 2.35\%，比独立假设下的 1.80\% 高了近 30\%。如果你忘了加交叉项，会用 3.25 的方差做决策——大幅低估了实际风险。

推广到任意有限个变量：

\[
\Var\!\left(\sum_{i=1}^n X_i\right)=\sum_i\Var(X_i)+2\sum_{i<j}\Cov(X_i,X_j).
\]

只有当所有变量两两不相关时，方差才退化为直接相加。\textbf{独立且二阶矩有限} 会推出协方差为零，所以独立变量的方差可加。但注意对比：\textbf{期望的可加性从来不要求独立}（见\hyperref[sec:05-Probability/04-Expectation-and-Moments/01]{``期望、线性性与指示变量''}）。期望求和什么都不管，方差求和则需要检查协方差。

若 \(X_1,\ldots,X_n\) 独立同分布，各方差 \(\sigma^2\)，样本均值 \(\bar X_n=\frac1n\sum X_i\) 的方差为

\[
\Var(\bar X_n)=\frac{\sigma^2}{n}.
\]

均值随样本量增大而稳定——但这以协方差消失为前提。如果观测共享一个系统偏差（比如全天的温度计都偏高 0.5 度），这个偏差在所有 \(X_i\) 中同时出现，协方差不为零，重复测量不会按 \(1/n\) 把它消除。是独立让 ``平均降噪'' 生效，不是平均本身。

\subsection{相关系数：消掉单位看强弱}

协方差有一个不直观的地方：它依赖单位。身高从米改为厘米，协方差瞬间膨胀 100 倍——但变量之间的关系强度没变。把协方差标准化，除以各自标准差，得到相关系数：

\[
\rho_{X,Y}=\frac{\Cov(X,Y)}{\sigma_X\sigma_Y},\qquad(\sigma_X,\sigma_Y>0).
\]

这是一个无量纲数。由 Cauchy--Schwarz 不等式（期望形式：\(|\E[UV]|\le\sqrt{\E[U^2]\E[V^2]}\)，取 \(U=X-\mu_X\)，\(V=Y-\mu_Y\)）：

\[
-1\le\rho_{X,Y}\le1.
\]

\(\rho=1\) 时二者完全正线性相关：存在 \(a>0,b\) 使 \(Y=aX+b\) 几乎必然成立。\(\rho=-1\) 时完全负线性相关。\(\rho=0\) 表示没有线性关系——仅此而已，不意味没有其他关系。

相关系数对平移和正缩放都不变：\(\rho_{aX+b,\,cY+d}=\rho_{X,Y}\)（\(a,c>0\)）。这使得它适合比较不同数据集、不同单位之间的关联强度。

\subsection{不相关远不等于独立}

反例是理解这句话最直接的方式。令 \(X\sim\mathrm{Unif}(-1,1)\)，定义 \(Y=X^2\)。显然 \(Y\) 完全由 \(X\) 决定——你知道 \(X\) 就知道 \(Y\) 精确值。强依赖。

但检查协方差：

\[
\E[X]=0,\qquad
\E[XY]=\E[X^3]=\int_{-1}^1\frac{x^3}{2}dx=0,
\]

奇函数在对称区间积分为零。所以

\[
\Cov(X,Y)=\E[XY]-\E[X]\E[Y]=0-0=0.
\]

协方差为零，相关系数也为零，可 \(Y\) 完全是 \(X\) 的函数。画散点图是一条抛物线，依赖关系清清楚楚，线性趋势却一丝不存。当你看到报告中 ``相关系数不足 0.1''，只能解读为 ``直线画不出什么趋势''，不能推断 ``两个量没关系''。

反向推呢？若 \((X,Y)\) 独立且各自的二阶矩有限，则

\[
\E[XY]=\E[X]\E[Y],
\]

（独立性的定义使联合求和拆成两个各自求和——与\hyperref[sec:05-Probability/04-Expectation-and-Moments/01]{``期望、线性性与指示变量''}中期望可加性不依赖独立形成了鲜明对照），因此 \(\Cov(X,Y)=0\)。独立推出不相关，但中间有一个特殊的例外族值得一提：若 \((X,Y)\) 服从联合 Gaussian 分布（联合密度完全由均值向量与协方差矩阵决定，详见\hyperref[sec:05-Probability/06-Common-Distributions/05]{``Gaussian 与多元 Gaussian''}），则在这一族内，不相关确实等价于独立。离开 Gaussian，再不能这样用。

\subsection{相关性可能被混合群体凭空制造}

某研究机构统计发现：街区内冰淇淋销量与溺水事故数正相关。难道吃冰淇淋导致溺水？更合理的解释是天气热的时候人们既买冰淇淋又去游泳——两个行为被第三个变量（气温）共同驱动。

在每个固定温度水平下，冰淇淋销量和溺水可能几乎不相关；一旦把不同温度的数据混在一起算总相关系数，就会出现虚假关联。极端时方向都可能反转：各年龄段内收入与健康指标正相关，但放进总人口后因为年龄同时反向影响两者，总体相关系数可能变负。这就是 Simpson 悖论的关联版本。

教训：相关系数依赖你选择的群体和纳入的条件信息。在作因果推断或解释相关性之前，至少应该问一句——有没有遗漏的共同原因在同时驱动两个变量？

下一节的任务是：如果你已经掌握了一些信息（比如温度、年龄），怎么在获取信息后仍然给出最佳的数值预测，并把总波动分解为 ``已知信息解释的部分'' 和 ``信息之外残余的部分''。
